\subsection{Human Expert Validation}\label{sec:human_expert_validation}

\paragraph{Quantitative survey statistics.}
\Cref{fig:expert_validation_results}
displays results from the human expert validations on \name{} (\texttt{gpt-oss-120b}) described in \Cref{sec:human_expert_validation_methods}.
Extraction Accuracy reports the average rate of field-level correctness conditional on the extraction being correctly flagged as relevant. In contrast, Expert-Rated Competence is assessed for every case, including incorrect flaggings.
\textit{Flagging Precision} reports the proportion of \name{} extractions judged relevant by experts and is higher for parameters ($0.66$) and outbreaks ($0.61$) than for models ($0.40$). \textit{Extraction Accuracy} reports the average rate of field-level correctness. All extraction types achieve expert-rated correctness comparable to or exceeding the precision of our automated Extraction evaluation in \Cref{sec:PERG_data_evaluation} ($0.77$ for parameters; $0.83$ for models; $0.80$ for outbreaks). Finally, average \textit{Expert-Rated Competence} is $4.2$ for parameters, $2.8$ for models, and $3.9$ for outbreaks.

\paragraph{Qualitative impressions.}
Based on survey feedback from six expert epidemiologists, \name{} was consistently reported to improve efficiency compared to fully manual extractions. Although false positives occur, these are typically easy to identify and remove, resulting in a net reduction in effort. Extraction difficulty varies across papers due to differences in complexity and reporting style, and in rare cases the system may increase effort in situations that are similarly challenging for human reviewers. Once an extraction is produced, individual fields are straightforward to validate, whereas multi-select fields pose more difficulty. Common errors arise from insufficient contextual information, limited use of document structure and cross-extraction constraints, and failures to infer fields apparent to human annotators when information is implicit. Additionally, the system struggles to understand provenance, occasionally mixing up newly reported findings with information cited from prior work.

% Expanded version of the above
%The system improves extraction efficiency compared to fully manual extractions. While it produces false positives, these are typically easy to identify and remove, preserving an overall net reduction in effort. Extraction difficulty varies by paper due to differences in complexity and reporting style, and in rare cases the system may increase effort in situations that are similarly non-trivial for human reviewers. An additional challenge arises from overlapping terminology within the field, including shared terms and units, which the system struggles to differentiate; extractions often appear to be driven by keyword presence rather than holistic reasoning.

%Once an extraction is produced, individual fields are straightforward to validate, whereas multi-select fields are more challenging. Errors arise from multiple sources. Insufficient global parameter context can cause fields to be populated in isolation without fully accounting for interdependencies, resulting in logically implausible extractions. Limited incorporation of document structure and weak enforcement of distinctness across extractions can lead to information being conflated across records. In addition, the system may fail to infer fields that are apparent to human annotators when such information is not stated explicitly in the text. Separately, weak restriction of provenance can cause newly reported findings to be mixed with information cited from prior work.


